\documentclass[../../main]{subfiles}

\begin{document}

\section{Relações Bem-Fundadas}
\label{sec:matematica:relacoes:relacoes-bem-fundadas}

\begin{defn}[Elemento Minimal]
  \label{def:matematica:relacoes:bem-fundadas:minimal}

  Seja (R) uma relação sobre (A).

  Um elemento (m\in A) é denominado minimal quando

  [
    \neg(\exists x\in A)(xRm).
  ]

\end{defn}

\begin{defn}[Relação Bem-Fundada]
  \label{def:matematica:relacoes:bem-fundadas}

  Uma relação (R) sobre um conjunto (A) é denominada bem-fundada
  quando todo subconjunto não vazio de (A) possui pelo menos um
  elemento minimal.

  Isto é,

  [
    (\forall B\subseteq A)
    \Bigl(
      B\neq\varnothing
      \Rightarrow
      (\exists m\in B)
      (\forall b\in B)
      \neg(bRm)
    \Bigr).
  ]

\end{defn}

\begin{prop}
  \label{prop:matematica:relacoes:bem-fundada-sem-cadeias}

  Se uma relação não admite cadeias descendentes infinitas, então ela é
  bem-fundada.

\end{prop}

\begin{obs}

  Em contextos clássicos com escolha dependente, a ausência de cadeias
  descendentes infinitas é equivalente à bem-fundação.

\end{obs}

% \begin{defn}[Boa Ordem]
%   \label{def:matematica:relacoes:boa-ordem}

%   Uma boa ordem é uma ordem total tal que todo subconjunto não vazio
%   possui um menor elemento.

% \end{defn}

% \begin{prop}
%   \label{prop:matematica:relacoes:boa-ordem-bem-fundada}

%   Toda boa ordem é uma relação bem-fundada.

% \end{prop}

% \begin{proof}

%   Seja (B\neq\varnothing).

%   Como a ordem é uma boa ordem, existe um menor elemento
%   (m\in B).

%   Em particular, não existe elemento de (B) estritamente menor que
%   (m).

%   Portanto (m) é minimal.

% \end{proof}

\begin{prop}[Princípio da Indução Bem-Fundada]
  \label{prop:matematica:relacoes:inducao-bem-fundada}

  Seja (R) uma relação bem-fundada sobre (A).

  Seja (P(x)) uma proposição definida para todo (x\in A).

  Se

  [
    (\forall x\in A)
    \Bigl(
      (\forall y\in A)
      (yRx \Rightarrow P(y))
      \Rightarrow
      P(x)
    \Bigr),
  ]

  então

  [
    (\forall x\in A)
    P(x).
  ]

\end{prop}

\begin{proof}

  Suponha que exista

  [
    a\in A
  ]

  tal que (P(a)) seja falsa.

  Considere o conjunto

  [
    S
    =
    {
      x\in A
      \mid
      \neg P(x)
    }.
  ]

  Como (R) é bem-fundada e (S\neq\varnothing), existe um elemento
  minimal (m\in S).

  Pela minimalidade de (m), todo elemento (y) com (yRm) satisfaz
  (P(y)).

  Pela hipótese indutiva, segue que (P(m)) é verdadeira, o que
  contradiz (m\in S).

  Portanto,

  [
    (\forall x\in A)P(x).
  ]

\end{proof}

\end{document}
